\metatitle{Newton polytopes and M-convexity}
\metadescription{Newton polytopes, saturated Newton polytopes, M-convex sets, generalized permutohedra, and their relation to Lorentzian polynomials.}


\section[newtonPolytopes]{Newton polytopes}

Given a polynomial
\[
  f=\sum_{\alpha} c_\alpha \xvec^\alpha \in \setC[x_1,\dotsc,x_n],
\]
with $\alpha\in\setN^n$, the \defin{support} of $f$ is
\[
  \operatorname{supp}(f)\coloneqq \{\alpha : c_\alpha \neq 0\}.
\]
The \defin{Newton polytope}\label{newtonPolytope} of $f$ is the convex hull
of the support.  We denote it by $\operatorname{Newton}(f)$.

Newton polytopes turn questions about monomial supports into questions about
lattice points and convex geometry.  In algebraic combinatorics, this is useful
for Schur, key, Schubert, Grothendieck, chromatic, and related polynomials.
The same language also appears naturally in
\hyperref[lorentzianPolynomial]{Lorentzian polynomials}, where the support is
controlled by \hyperref[MConvexity]{$M$-convexity}.


\section[SNP]{Saturated Newton polytopes}

A polynomial has a \defin{saturated Newton polytope} (SNP) if
\[
  \operatorname{Newton}(f)\cap \setZ^n = \operatorname{supp}(f).
\]
Equivalently, every lattice point in the Newton polytope occurs as the exponent
vector of a monomial with nonzero coefficient.

Several classical polynomials, such as \hyperref[schurS]{Schur polynomials} and
\hyperref[stanleySym]{Stanley symmetric functions}, have saturated Newton
polytopes.  A foundational source for this circle of questions is the paper of
\name{Cara Monical}, \name{Neriman Tokcan}, and \name{Alexander Yong}
\cite{MonicalTokcanYong2019}, where many conjectures were formulated.

Key polynomials and Schubert polynomials have SNP
\cite{FinkMeszarosDizier2018}.  Key polynomials in other types
($A_r$, $B_r$, $C_r$, $F_4$, $G_2$) also have SNP; see
\cite{BessonJeraldsKiers2023}.
\name{Greta Panova} and \name{Chenchen Zhao} prove special cases of the
Monical--Tokcan--Yong conjecture that Kronecker products of Schur functions
have saturated Newton polytopes \cite{PanovaZhao2023x}.  Their arguments use
Horn inequalities for Littlewood--Richardson coefficients and imply necessary
conditions for positivity of Kronecker coefficients.

\name{Mario Sanchez} studies when matroid polytopes, flag matroid polytopes,
Bruhat interval polytopes, and Schubitopes are Mirković--Vilonen polytopes
\cite{Sanchez2023x}.  In particular, the Newton polytopes of Schubert
polynomials and key polynomials are Mirković--Vilonen polytopes.

\name{Marc Besson}, \name{Sam Jeralds}, and \name{Joshua Kiers} extend part
of this Demazure-polytope picture to affine Demazure modules
\cite{BessonJeraldsKiers2024x}.  Their affine Demazure weight polytopes are
cut out by inequalities governed by standard, opposite, and semi-infinite
Bruhat orders, while their face vertices are described using twisted Bruhat
orders.

SNP for \hyperref[chromaticQuasisymmetric]{chromatic symmetric polynomials}
obtained from incomparability graphs of $(3+1)$-free posets was proved by
\name{Jacob P. Matherne}, \name{Alejandro H. Morales}, and
\name{Jesse Selover} \cite{MatherneMoralesSelover2022x}.

For more background and results, see \cite{WangZhangZhang2024x}.
In \cite{NguyenNgocTuanHai2023}, the authors prove SNP for many families of
symmetric functions, such as \hyperref[schurP]{Schur-$P$} and
\hyperref[schurQ]{Schur-$Q$} polynomials, as well as many of the families
above.  Moreover, they prove that $\operatorname{Newton}(f)$ has the
\hyperref[polytopeIDP]{integer decomposition property} in many cases.


\section[MConvexity]{$M$-convexity}

A subset $J\subseteq\setN^n$ is \defin{$M$-convex}\label{mConvexSet} if for
all $\alpha,\beta\in J$, and any index $i\in[n]$ with
$\alpha_i>\beta_i$, there is some index $j$ with $\alpha_j<\beta_j$ and
\[
  \alpha-e_i+e_j\in J.
\]
Here $e_i$ is the $i^\thsup$ unit vector.  This is the integer-lattice version
of the exchange axiom: for a \hyperref[matroid]{matroid}, the indicator
vectors of the bases form an $M$-convex subset of $\{0,1\}^n$.

\begin{theorem}[See \cite{DizierPhDThesis}]
A homogeneous polynomial $f$ is $M$-convex if and only if $f$ has a saturated
Newton polytope and this Newton polytope is a
\hyperref[generalizedPermutohedron]{generalized permutohedron}.
\end{theorem}

\hyperref[key]{Key polynomials} are $M$-convex; this was conjectured earlier by
\name{Cara Monical}, \name{Neriman Tokcan}, and \name{Alexander Yong}
\cite{MonicalTokcanYong2019}, and proved in \cite{FanGuoPengSun2020}.
In \cite{WangZhangZhang2024x}, it is proved that affine Stanley symmetric
functions and cylindric skew Schur functions are $M$-convex.

\name{Serena An}, \name{Katherine Tung}, and \name{Yuchong Zhang}
characterize the support of dual Schubert polynomials as a Minkowski sum over
the inversions of the indexing permutation \cite[Thm.~1.2]{AnTungZhang2024x}.
This gives an elementary proof that dual Schubert polynomials are
$M$-convex, have saturated Newton polytopes, and have Newton polytopes which
are generalized permutohedra \cite[Cor.~1.3]{AnTungZhang2024x}.  They also
give a combinatorial description of the vertices of these Newton polytopes.


\section[newtonLorentzianBridge]{Relation with Lorentzian polynomials}

The support condition in the definition of a
\hyperref[lorentzianPolynomial]{Lorentzian polynomial} is exactly
$M$-convexity.  Thus Lorentzian polynomials combine a discrete convex support
condition with a Hessian-signature condition.  In particular, if $B$ is the set
of bases of an integral \hyperref[polymatroids]{polymatroid}, then the
normalized support polynomial
\[
  \sum_{\alpha\in B}
  \frac{\xvec^\alpha}{\alpha_1!\alpha_2!\dotsm\alpha_n!}
\]
is Lorentzian \cite{BrandenHuh2020}.  For ordinary matroids this recovers the
\hyperref[basisGeneratingPolynomial]{basis-generating polynomial}.
